\begin{enumerate}
\item Write an inductive proof of the formula for the sum 
of the first $n$ cubes.

\hint{
\begin{thm*}
\[ \forall n \in \Naturals, \; \sum_{k=1}^n k^3 \;= \; \left( \frac{n(n+1)}{2} \right)^2 \]
\end{thm*}

\begin{proof} (By mathematical induction)

{\bf Base case:} ($n=1$)
For the base case, note that when $n=1$ we have

\[ \sum_{k=1}^n k^3 \; = \; 1 \]

\noindent and 

\[ \left( \frac{n(n+1)}{2} \right)^2  \; = \; 1.\]

{\bf Inductive step:}

Suppose that $m>1$ is an integer such that 

\[ \sum_{k=1}^m k^3 \; = \; \left( \frac{m(m+1)}{2} \right)^2 \]

\noindent Add $(m+1)^3$ to both sides to obtain

\[ (m+1)^3 + \sum_{k=1}^m k^3 \;= \; \left( \frac{m(m+1)}{2} \right)^2 + (m+1)^3. \]

\noindent Thus

\begin{gather*} 
\sum_{k=1}^{m+1} k^3 \;= \; \left( \frac{m^2(m+1)^2}{4} \right) + \frac{4(m+1)^3}{4} \\
\;= \; \left( \frac{m^2(m+1)^2 + 4 (m+1)^3}{4} \right)\\
\; = \; \left( \frac{(m+1)^2 (m^2 + 4(m+1))}{4} \right)\\
\; = \; \left( \frac{(m+1)^2 (m^2 + 4m +4)}{4} \right)\\
\; = \; \left( \frac{(m+1)^2 (m+2)^2}{4} \right)\\
\; = \; \left( \frac{(m+1)(m+2)}{2} \right)^2
\end{gather*}

 \end{proof}
}

\wbvfill

\workbookpagebreak
  
\item Find a formula for the sum of the first $n$ fourth powers.

\hint{\[ \frac{n\cdot(n+1)\cdot(2n+1)\cdot(3n^2+3n-1)}{30}\] } 

\wbvfill

\workbookpagebreak
 
\item The sum of the first $n$ natural numbers is sometimes called
the $n$-th triangular number \index{triangular numbers}$T_n$.  Triangular numbers are so-named
because one can represent them with triangular shaped arrangements 
of dots. 

\begin{center} \begin{picture}(0,0)%
\includegraphics{./triangular_numbers.pdf}%
\end{picture}%
\setlength{\unitlength}{3947sp}%
%
\begingroup\makeatletter\ifx\SetFigFont\undefined%
\gdef\SetFigFont#1#2#3#4#5{%
  \reset@font\fontsize{#1}{#2pt}%
  \fontfamily{#3}\fontseries{#4}\fontshape{#5}%
  \selectfont}%
\fi\endgroup%
\begin{picture}(5963,1088)(1193,-2116)
\end{picture}%
 \end{center}

The first several triangular numbers are 1, 3, 6, 10, 15, et cetera.

Determine a formula for the sum of the first $n$ triangular numbers $\displaystyle \left( \sum_{i=1}^n T_i \right)$ and prove it using PMI.

\hint{The formula is $\frac{n(n+1)(n+2)}{6}$.}

\wbvfill

\workbookpagebreak

\item Consider the alternating sum of squares:
\begin{gather*}
1 \\
1 - 4 = -3 \\
1 - 4 + 9 = 6 \\
1 - 4 + 9 - 16 = -10 \\
\mbox{et cetera}
\end{gather*}
Guess a general formula for $\sum_{i=1}^n (-1)^{i-1} i^2$, and prove it using PMI.

\hint{
\begin{thm*}
\[ \forall n \in \Naturals, \; \sum_{i=1}^n (-1)^{i-1} i^2 \;= \; (-1)^{n-1} \frac{n(n+1)}{2}  \]
\end{thm*}

\begin{proof} (By mathematical induction)

{\bf Base case:} ($n=1$)
For the base case, note that when $n=1$ we have

\[\sum_{i=1}^n (-1)^{i-1} i^2 \;= \; 1 \]

\noindent and also

\[ (-1)^{n-1} \frac{n(n+1)}{2} \;= \; 1. \]

{\bf Inductive step:}

Suppose that $k>1$ is an integer such that 

\[ \sum_{i=1}^k (-1)^{i-1} i^2 \;= \; (-1)^{k-1} \frac{k(k+1)}{2}.  \]

Adding $(-1)^{k} (k+1)^2$ to both sides gives

\begin{gather*} 
\sum_{i=1}^{k+1} (-1)^{i-1} i^2 \;= \; (-1)^{k-1} \frac{k(k+1)}{2} + (-1)^{k} (k+1)^2 \\
\;= \; (-1)^{k-1} \frac{k(k+1)}{2} - (-1)^{k-1} (k+1)^2 \\ 
\;= \; (-1)^{k-1} \left( \frac{k(k+1)}{2} -  \frac{2(k+1)^2}{2} \right) \\ 
\;= \; (-1)^{k} \left( \frac{2(k+1)^2}{2} - \frac{k(k+1)}{2} \right) \\
\;= \; (-1)^{k} \frac{(k+1)(2(k+1)-k)}{2} \\
\;= \; (-1)^{k} \frac{(k+1)(k+2)}{2} \\
\end{gather*}
\end{proof}
}

\wbvfill

\workbookpagebreak

\item Prove the following formula for a product.

\[ \prod_{i=2}^n \left(1 - \frac{1}{i}\right) =  \frac{1}{n} \]

\wbvfill


\item Prove $\displaystyle \sum_{j=0}^{n}(4j+1) \; = \; 2n^{2}+3n+1$ for all
integers $n \geq 0$.

\wbvfill

\workbookpagebreak

\item Prove $\displaystyle \sum_{i=1}^{n}\frac{1}{(2i-1)(2i+1)} \; = \; \frac{n}{2n+1}$ for all natural numbers $n$.

\wbvfill

\workbookpagebreak

\item The \index{Fibonacci numbers} \emph{Fibonacci numbers} are a sequence of integers defined by
the rule that a number in the sequence is the sum of the two that 
precede it.

\[ F_{n+2} = F_n + F_{n+1}  \]

\noindent The first two Fibonacci numbers (actually the zeroth and the first) 
are both 1.  

\noindent Thus, the first several Fibonacci numbers are

\[ F_0 = 1, F_1=1, F_2=2, F_3=3, F_4=5, F_5=8, F_6=13, F_7=21, \; \mbox{et cetera} \]

Use mathematical induction to prove the following formula involving
Fibonacci numbers.

\[ \sum_{i=0}^n (F_i)^2 \, = \, F_n \cdot F_{n+1} \]

\wbvfill

\workbookpagebreak

\end{enumerate}

%% Emacs customization
%% 
%% Local Variables: ***
%% TeX-master: "GIAM-hw.tex" ***
%% comment-column:0 ***
%% comment-start: "%% "  ***
%% comment-end:"***" ***
%% End: ***

